% ─────────────────────────────────────────────────────────────
% §Background — ONE subsection naming the common flaw shared by Paradigm I (monolithic bloat)
% and Paradigm II (flat-bank mis-routing): both lack governance. ~80 words.
% \input right after \section{Introduction}. The named concept "Skill Bloat vs. Skill Governance"
% is introduced in Intro 段2 句2.6 and crystallized here as a standalone subsection (brief Part 0.6).
% ─────────────────────────────────────────────────────────────

\subsection{Skill Bloat vs.\ Skill Governance}
\label{sec:background}
Both dominant paradigms share one root flaw: they treat skill evolution as ungoverned \emph{accumulation}.
\textbf{(I) Monolithic} methods merge every distilled patch into a single ever-growing document, so growth becomes \textbf{skill bloat}---irrelevant procedure is loaded in full and drives negative transfer.
\textbf{(II) Flat skill banks} store skills as independent entries and retrieve a top-$k$ slice, so growth becomes \textbf{mis-routing}---with no dependency closure and no lifecycle, similarity-based selection omits prerequisites and admits conflicts.
Neither can split a coarse skill, route its dependencies, or retire a harmful one; we call this missing capability \textbf{skill governance}.
